\section{Equilibria without an efficiency bound}
\label{app:rewrite-uncapped-proofs}
\Needspace{7\baselineskip}
\begin{proof}[Proof of Proposition~\ref{prop:rewrite-uncapped-complements-bounds}]
\phantomsection
\label{proof:rewrite-uncapped-complements-bounds}
Dropping the positive AI term from the CES sum gives an upper bound for
$Z$, because the outer exponent $\sigma/(\sigma-1)$ is negative.
This gives the production bounds in
\eqref{eq:rewrite-uncapped-complements-z-bound}
and~\eqref{eq:rewrite-uncapped-complements-y-bound}.
Define $k=K/(AL)$. The resource constraint
\eqref{eq:rewrite-eq-kdot}, nonnegative consumption and compute expenditure,
and $L=N$ imply
\begin{equation}
 \dot k\leq
 \omega_L^{\sigma(1-\alpha)/(\sigma-1)}k^\alpha
 -(\delta+n+\gamma)k.
 \label{eq:rewrite-uncapped-complements-k-bound}
\end{equation}
Since $\delta+n+\gamma>0$ and $0<\alpha<1$, the right-hand side is
negative whenever $k$ exceeds a finite positive level. In particular,
\[
 k_t\leq
 \max\left\{k_0,
 \left[\frac{\omega_L^{\sigma(1-\alpha)/(\sigma-1)}}
 {\delta+n+\gamma}\right]^{1/(1-\alpha)}\right\}.
\]
Equation~\eqref{eq:rewrite-uncapped-complements-y-bound} then bounds
$Y/(AL)$ uniformly. Since $A_t=A_0e^{\gamma t}$ and $L=N$,
$Y/N$ is bounded above by a constant times $e^{\gamma t}$.
Taking logarithms, dividing by $t$, and taking the upper limit gives
\eqref{eq:rewrite-uncapped-complements-growth-bound}.
\end{proof}

% Additional results retained for possible later use; not part of the proposition.
% The same bounds keep $K$ and $Y$ finite on every bounded time interval.
% For any such horizon $T$, integrating the resource constraint gives
% \[
%  \int_0^T M_t\dd t\leq K_0+\int_0^T Y_t\dd t<\infty.
% \]
% The uncapped RSI technology and H\"older's inequality imply
% \begin{align}
%  B_T^{1-\eta}
%  &=B_0^{1-\eta}+(1-\eta)\chi\int_0^T M_t^\eta\dd t
%  \nonumber\\
%  &\leq B_0^{1-\eta}
%  +(1-\eta)\chi T^{1-\eta}
%  \left(\int_0^T M_t\dd t\right)^\eta<\infty.
%  \label{eq:rewrite-uncapped-complements-b-bound}
% \end{align}
% Thus $B$ cannot diverge at a finite date either. This argument bounds
% stocks and output; it is not a regularity theorem for all instantaneous
% control choices or an equilibrium-existence proof.
%
% There is also no finite-horizon unbounded-payoff argument of the kind used
% for substitution in Proposition~\ref{prop:rewrite-research-scale}.
% At given positive continuous paths of $K,A,L$ on $[0,T]$, optimized operating
% profit satisfies
% \[
%  0\leq\pi_t(B)\leq p_XX\leq(1-\alpha)Y
%  \leq(1-\alpha)\omega_L^{\sigma(1-\alpha)/(\sigma-1)}
%  K_t^\alpha(A_tL_t)^{1-\alpha}.
% \]
% This upper bound is independent of the developer's choice of $B$.
% For a given locally integrable $r$, the discount factor
% $D_t=\exp(-\int_0^t r_s\dd s)$ has a positive minimum on $[0,T]$.
% Discounted operating revenue is bounded, whereas the cost of research
% tends to infinity with $\int_0^T M_t\dd t$. The developer's truncated
% supremum is therefore finite, and net payoff tends to minus infinity as
% total research expenditure diverges. This conclusion holds for every
% $0<\eta<1$, but does not prove that an infinite-horizon optimum exists.

\Needspace{7\baselineskip}
\begin{proof}[Proof of Proposition~\ref{prop:rewrite-uncapped-unit-bgp}]
\phantomsection
\label{proof:rewrite-uncapped-unit-bgp}
I first construct and verify the reference BGP, denoted by a superscript
$*$, and then establish convergence from nearby initial stocks.

\textit{Step 1. Construction of the reference BGP.}
At unit elasticity, $s_X=\omega_X$ and $e_X=1-(1-\alpha)\omega_X$.
The AI-demand and monopoly conditions,
equations~\eqref{eq:rewrite-ai-price} and~\eqref{eq:rewrite-monopoly-foc}, give
\begin{equation}
 \frac{p_XX}{Y}=(1-\alpha)\omega_X,\qquad
 \frac UY=(1-\alpha)^2\omega_X^2.
 \label{eq:rewrite-uncapped-unit-static-shares}
\end{equation}
Substituting $X=BU=(1-\alpha)^2\omega_X^2BY$ into production yields
\begin{equation}
 Y^{1-(1-\alpha)\omega_X}
 =K^\alpha(AL)^{(1-\alpha)\omega_L}
 \bigl[(1-\alpha)^2\omega_X^2B\bigr]^{(1-\alpha)\omega_X}.
 \label{eq:rewrite-uncapped-unit-production}
\end{equation}
On an exact BGP with positive constant $K/Y$, $C/Y$, and $M/Y$,
log differentiation of this equation and the RSI technology
\eqref{eq:rewrite-uncapped-unit-law} gives the growth restrictions
\eqref{eq:rewrite-uncapped-unit-production-growth}
and~\eqref{eq:rewrite-uncapped-unit-research-growth}.
Define
\begin{equation}
 \Delta\equiv(1-\alpha)(1-\eta-\omega_X)>0.
 \label{eq:rewrite-uncapped-unit-feedback}
\end{equation}
The unique solution to the two growth restrictions has
\begin{equation}
 g_Y^*=\frac{(1-\eta)(1-\omega_X)}{1-\eta-\omega_X}(n+\gamma),
 \label{eq:rewrite-uncapped-unit-output-growth}
\end{equation}
and $g_B^*$ as in \eqref{eq:rewrite-uncapped-unit-efficiency-growth}.
The equality $g_K=g_C=g_Y^*$ follows from constant capital and consumption ratios;
equations~\eqref{eq:rewrite-wage} and~\eqref{eq:rewrite-euler} then give
\eqref{eq:rewrite-uncapped-unit-wage-growth}
and~\eqref{eq:rewrite-uncapped-unit-return}.
Since $n+\gamma>0$, the two growth restrictions are inconsistent if
$\Delta=0$ and imply $g_B<0$ if $\Delta<0$, contrary to the RSI technology.
Thus $\eta+\omega_X\geq1$ excludes this positive finite-rate BGP,
not every possible equilibrium trajectory.

Define the candidate ratios $k^*=K/Y$, $u^*=U/Y$,
$i^*=(\dot K+\delta K)/Y$, $m^*=M/Y$, and $c^*=C/Y$ by
\begin{align}
 k^*&=\frac{\alpha}{r^*+\delta},\qquad u^*=(1-\alpha)^2\omega_X^2,
 \label{eq:rewrite-uncapped-unit-capital-inference}\\
 i^*&=(g_Y^*+\delta)k^*,
 \label{eq:rewrite-uncapped-unit-investment}\\
 m^*&=\frac{u^*\eta g_B^*}{\rho-n+\eta g_Y^*},
 \label{eq:rewrite-uncapped-unit-research-share}\\
 c^*&=1-u^*-i^*-m^*.
 \label{eq:rewrite-uncapped-unit-consumption}
\end{align}
Here $k^*$ is a capital--output ratio, not capital per unit of effective
labor, and $i^*$ is gross physical investment as a share of output.
The research share comes from the developer's first-order and costate
conditions; the consumption share comes from market clearing.
These ratios are endogenous, not fixed saving or research rules.

All candidate ratios are positive. The denominator of the research share $m^*$,
$\rho-n+\eta g_Y^*$, is positive. Moreover,
\[
 i^*=\frac{\alpha(g_Y^*+\delta)}{\rho-n+g_Y^*+\delta}<\alpha,
 \qquad
 m^*<\frac{u^*\eta}{1-\eta}.
\]
Since $\Delta=(1-\alpha)(1-\eta-\omega_X)>0$,
$\omega_X<1-\eta$ and $\omega_X<1$. Consequently,
\[
 u^*+m^*<\frac{u^*}{1-\eta}<(1-\alpha)\omega_X<1-\alpha,
 \qquad c^*=1-u^*-i^*-m^*>0.
\]
Thus positive consumption follows from the stated conditions; it is not
an additional assumption.

For given $A_0,N_0>0$, choose the reference initial stocks as
\begin{align}
 B_0^*&=\left[
 A_0N_0(k^*)^{\alpha/[(1-\alpha)\omega_L]}(u^*)^{\omega_X/\omega_L}
 m^*\left(\frac{\chi}{g_B^*}\right)^{1/\eta}
 \right]^{\eta\omega_L/(1-\eta-\omega_X)},
 \label{eq:rewrite-uncapped-unit-initial-b}\\
 Y_0^*&=A_0N_0(k^*)^{\alpha/[(1-\alpha)\omega_L]}
 (u^*B_0^*)^{\omega_X/\omega_L},
 \qquad K_0^*=k^*Y_0^*.
 \label{eq:rewrite-uncapped-unit-initial-y-k}
\end{align}
The outer exponent in the expression for $B_0^*$ is positive. These
choices solve simultaneously the initial production equation and
$g_B^*B_0^*=\chi(B_0^*m^*Y_0^*)^\eta$.
Set $U_0^*=u^*Y_0^*$, $M_0^*=m^*Y_0^*$, $C_0^*=c^*Y_0^*$,
$X_0^*=B_0^*U_0^*$, and
\begin{equation}
 q_0^*=\frac{M_0^*}{\eta g_B^*B_0^*}.
 \label{eq:rewrite-uncapped-unit-initial-q}
\end{equation}
Let $Y_t^*=Y_0^*e^{g_Y^*t}$, $B_t^*=B_0^*e^{g_B^*t}$,
$K_t^*=k^*Y_t^*$, $C_t^*=c^*Y_t^*$, $U_t^*=u^*Y_t^*$,
$M_t^*=m^*Y_t^*$, and $q_t^*=q_0^*e^{(g_Y^*-g_B^*)t}$.
Use the exogenous $A_t,N_t$ and set $L_t=N_t$ and $X_t^*=B_t^*U_t^*$.
The static equilibrium equations determine all remaining quantities and prices.

The two growth identities
\eqref{eq:rewrite-uncapped-unit-production-growth}
and~\eqref{eq:rewrite-uncapped-unit-research-growth} ensure that the production
equation and the RSI technology hold at every date once they hold initially.
Equation~\eqref{eq:rewrite-uncapped-unit-initial-q} gives
$M=\eta q\dot B$, which is the research first-order condition
\eqref{eq:rewrite-eq-research} with $\psi=1$.
The definitions of $k^*$, $r^*$, and $c^*$ verify the capital-return,
Euler, and resource equations. The costate equation reduces to
\[
 g_Y^*-g_B^*
 =r^*-\frac{u^*\eta g_B^*}{m^*}-\eta g_B^*.
\]
This holds by equation~\eqref{eq:rewrite-uncapped-unit-research-share}
and $(1-\eta)g_B^*=\eta g_Y^*$.
All allocations are positive and finite at every finite date.

Both transversality conditions are explicit:
\begin{align}
 e^{-\rho t}\frac{N_tK_t^*}{C_t^*}
 &=\frac{N_0k^*}{c^*}e^{-(\rho-n)t}\longrightarrow0,
 \label{eq:rewrite-uncapped-unit-household-tvc}\\
 e^{-r^*t}q_t^*B_t^*
 &=q_0^*B_0^*e^{-(\rho-n)t}\longrightarrow0.
 \label{eq:rewrite-uncapped-unit-developer-tvc}
\end{align}
Discounted developer profit is a constant multiple of
$e^{-(\rho-n)t}$. Discounted household utility is
$N_0e^{-(\rho-n)t}[\log(C_0^*/N_0)+(g_Y^*-n)t]$.
Both improper objective integrals are finite.

\par\medskip\noindent\textit{Step 2. Global optimality.}

The constructed path satisfies the necessary conditions, but these do
not yet rule out profitable deviations. For the developer, holding the
paths of $K,A,L$ fixed, operating profit is AI revenue less inference
expenditure. Substituting the static shares
\eqref{eq:rewrite-uncapped-unit-static-shares} and then eliminating output
using~\eqref{eq:rewrite-uncapped-unit-production} gives
\begin{equation}
 \begin{split}
 \pi_t(B)={}&(1-\alpha)\omega_X[1-(1-\alpha)\omega_X]
 \bigl[K_t^\alpha(A_tL_t)^{(1-\alpha)\omega_L}\bigr]^{1/[1-(1-\alpha)\omega_X]}\\
 &\times\bigl[(1-\alpha)^2\omega_X^2B\bigr]^{(1-\alpha)\omega_X/[1-(1-\alpha)\omega_X]}.
 \end{split}
 \label{eq:rewrite-uncapped-unit-profit}
\end{equation}
To verify developer optimality by concavity, define the invertible state
transformation and its current-value costate by
\begin{equation}
 \nu=\frac{1-\eta}{\eta},\qquad
 S=B^\nu,\qquad p=\frac{q}{\nu B^{\nu-1}}.
 \label{eq:rewrite-uncapped-unit-transformed-state}
\end{equation}
The exponent $\nu$ is chosen so that the transformed research technology
has constant returns in $S$ and $M$:
\[
 \dot S=\nu\chi S^{1-\eta}M^\eta,
\]
which is jointly concave in $S,M$ for every $0<\eta<1$.
At given $K,A,L$, operating profit in
equation~\eqref{eq:rewrite-uncapped-unit-profit} is a positive
time-dependent coefficient times
\[
 S^{\frac{(1-\alpha)\omega_X}
 {\nu[1-(1-\alpha)\omega_X]}}.
\]
Since $\eta+\omega_X<1$ implies $(1-\alpha)\omega_X<1-\eta$, this
exponent lies strictly between zero and one. Operating profit is therefore
concave in $S$. With $p>0$, the transformed Hamiltonian
$\pi_t(S^{1/\nu})-M+p_t\nu\chi S^{1-\eta}M^\eta$ is jointly concave in $S,M$.
The research condition~\eqref{eq:rewrite-research-foc} and costate
equation~\eqref{eq:rewrite-costate}, with $\psi=1$, give the transformed
problem's control first-order condition and costate equation by the chain rule.

Write $D_t=\exp(-\int_0^t r_s\dd s)$ and
$J_T=\int_0^T D_t[\pi_t(B_t)-M_t]\dd t$.
For any admissible alternative with the same initial $B$, concavity,
the research condition, and the costate equation give, after integration
by parts,
\begin{equation}
 \widetilde J_T-J_T
 \leq-D_Tp_T(\widetilde S_T-S_T)
 \leq D_Tp_TS_T
 =\frac{D_Tq_TB_T}{\nu}\longrightarrow0.
 \label{eq:rewrite-uncapped-unit-developer-verification}
\end{equation}
The last limit follows from the developer TVC. Since the transformation
preserves feasible research paths and their payoffs, no admissible
alternative has a greater limiting payoff, or a greater upper limit of
truncated payoff, than this finite-valued candidate. The same verification
applies to nearby candidates satisfying the uncapped equations and the
developer TVC; it does not require exact balanced growth.

Household optimality follows from log-utility concavity, the linear budget,
the Euler equation, and the TVC, by the same integration-by-parts argument
as in Step 4 of the proof of Proposition~\ref{prop:rewrite-equilibrium-regimes}.

Competitive final-good firms
face a concave constant-returns technology, and
Lemma~\ref{lem:rewrite-monopoly-choice} establishes global static monopoly
optimality. The constructed path satisfies Definition~\ref{def:rewrite-equilibrium}
with the uncapped RSI technology~\eqref{eq:rewrite-uncapped-unit-law}.

\medskip
\textit{Step 3. Local stability and selection of initial consumption and research.}
\phantomsection
\label{proof:rewrite-uncapped-unit-local}
The initial stocks are predetermined. I must select initial consumption
and the shadow price so that the resulting path converges to the BGP;
the research condition then determines initial research expenditure.
For the given exogenous paths of population and labor-augmenting technology,
compare each variable with its value on the reference BGP at the same date.
Define log deviations
\begin{equation}
 \xi_K=\log\frac K{K^*},\quad
 \xi_B=\log\frac B{B^*},\quad
 \xi_C=\log\frac C{C^*},\quad
 \xi_q=\log\frac q{q^*}.
 \label{eq:rewrite-uncapped-unit-deviations}
\end{equation}
Time subscripts are suppressed: the reference quantities in the denominators
follow the BGP rather than remain fixed at their initial values.
Substitution into equations~\eqref{eq:rewrite-eq-kdot}--\eqref{eq:rewrite-eq-qdot},
with $\psi=1$ and $\psi'=0$, gives a smooth autonomous system
$\dot\xi=f(\xi)$ with $f(0)=0$, where $\xi=(\xi_K,\xi_B,\xi_C,\xi_q)'$.
Let $J^*=Df(0)$ denote its Jacobian. No constant growth rates or
allocation shares are imposed along these transitions.

I next establish two properties needed for this selection. First, the
linearized system has exactly two independent directions of convergence
and no eigenvalues with zero real part. Writing $E^s$ for the real
invariant subspace associated with the convergent directions, this means
\begin{equation}
 \operatorname{spec}(J^*)\cap i\mathbb R=\varnothing,
 \qquad \dim E^s=2.
 \label{eq:rewrite-uncapped-unit-local-spectrum}
\end{equation}
Second, let $V_S\in\mathbb R^{4\times2}$ be any real basis matrix for
$E^s$, and let $V_{KB}$ consist of its first two rows. I show that
\begin{equation}
 \det V_{KB}\ne0.
 \label{eq:rewrite-uncapped-unit-local-projection}
\end{equation}
This nonsingularity allows the stable manifold to select $C_0,q_0$ for
every sufficiently close pair $K_0,B_0$, rather than only for stocks on a
curve. Both properties follow from the parameter restrictions, rather
than being additional assumptions. I establish them by locating the
stable eigenvalues from the signs of the characteristic polynomial,
then checking that their eigenvectors allow the two stocks to vary
independently. Retain $\nu$ from
\eqref{eq:rewrite-uncapped-unit-transformed-state} and, for this linearization,
define
\begin{align*}
 a_K&=\frac{\alpha}{1-(1-\alpha)\omega_X},&
 a_T&=\frac{(1-\alpha)\omega_X\eta}
 {[1-(1-\alpha)\omega_X](1-\eta)},\\
 d&=\rho-n,&P&=r^*+\delta,&h&=\rho-n+\eta g_Y^*.
\end{align*}
The quantities $\nu,a_K,a_T,d,P,h$ are positive, and the restriction
$\eta+\omega_X<1$ gives
\begin{equation}
 1-a_K-a_T
 =\frac{(1-\alpha)(1-\eta-\omega_X)}
 {[1-(1-\alpha)\omega_X](1-\eta)}>0.
 \label{eq:rewrite-uncapped-unit-local-gap}
\end{equation}
For this calculation, return to the transformed state $S$ and its shadow
price $p$ from Step 2. This gives the invertible linear change of deviations
\[
 \widehat\xi=(\xi_K,\nu\xi_B,\xi_C,\xi_q+(1-\nu)\xi_B)'.
\]
The second and fourth coordinates are the log deviations of the transformed
state $S$ and its shadow price $p$. This changes neither the model nor the
eigenvalues. In these coordinates, the reduced production equation and
the research condition give
\[
 \log\frac Y{Y^*}=a_K\widehat\xi_K+a_T\widehat\xi_S,\qquad
 \log\frac M{M^*}=\widehat\xi_S+\frac{\widehat\xi_p}{1-\eta}.
\]
Differentiating the capital, transformed research, Euler, and transformed
costate equations at the BGP, using
$U/(qB)=\rho-n+\eta g_Y^*$, gives a matrix similar to $J^*$:
\begin{equation}
 \widehat J=
 \begin{pmatrix}
  \mathsf A&\mathsf B&-\mathsf C&-\mathsf V\\
  0&0&0&\mathsf D\\
  -\mathsf G&\mathsf H&0&0\\
  -\mathsf E&\mathsf F&0&d
 \end{pmatrix},
 \label{eq:rewrite-uncapped-unit-local-similarity}
\end{equation}
where the following letters abbreviate matrix coefficients only:
\begin{align*}
 \mathsf A&=d+(1-\alpha)\omega_X P,&
 \mathsf B&=\frac{(1-u^*)a_T-m^*}{k^*},&
 \mathsf C&=\frac{c^*}{k^*},\\
 \mathsf V&=\frac{m^*}{k^*(1-\eta)},&
 \mathsf D&=\frac{\eta g_Y^*}{1-\eta},&
 \mathsf G&=P(1-a_K),\\
 \mathsf H&=Pa_T,&
 \mathsf E&=\mathsf G+ha_K,&
 \mathsf F&=\mathsf H+h(1-a_T).
\end{align*}
All coefficients other than $\mathsf B$ are positive by their definitions.
The BGP allocation
ratios imply the identities
\begin{align}
 \mathsf L&\equiv\mathsf B-\frac{\mathsf Vd}{\mathsf D}
 =\frac{(1-\alpha)\omega_X\eta}{(1-\eta)k^*}>0,
 \label{eq:rewrite-uncapped-unit-local-l}\\
 \mathsf G\mathsf F-\mathsf H\mathsf E
 &=hP(1-a_K-a_T)>0.
 \label{eq:rewrite-uncapped-unit-local-minor}
\end{align}
Set $\mathsf N=\mathsf C\mathsf G+\mathsf V\mathsf E>0$ and
$\mathsf Q=\mathsf C\mathsf H+\mathsf V\mathsf F>0$.
Expanding the determinant gives the characteristic polynomial
\begin{equation}
 \begin{split}
 p(\ell)&=\det(\ell I-\widehat J)\\
 &=(\ell^2-\mathsf A\ell-\mathsf N)
 (\ell^2-d\ell-\mathsf D\mathsf F)
 -\mathsf D\mathsf E(\mathsf Q-\mathsf L\ell).
 \end{split}
 \label{eq:rewrite-uncapped-unit-local-polynomial}
\end{equation}
To locate the negative eigenvalues without solving the quartic, first
evaluate the polynomial at zero:
\[
 p(0)=\mathsf D\mathsf C hP(1-a_K-a_T)>0.
\]
Let $\mu=(d-\sqrt{d^2+4\mathsf D\mathsf F})/2<0$, the negative
root of $\ell^2-d\ell-\mathsf D\mathsf F$. At $\ell=\mu$, the product
of the two quadratic factors in $p(\ell)$ vanishes, leaving
$p(\mu)=-\mathsf D\mathsf E(\mathsf Q-\mathsf L\mu)<0$.
Since $p(\ell)\to+\infty$ as $\ell\to-\infty$, there is one
negative root below $\mu$ and another between $\mu$ and zero.
These are the only negative roots: the cubic coefficient of $p$ is
$-(\mathsf A+d)<0$, its linear coefficient is
$\mathsf A\mathsf D\mathsf F+d\mathsf N+\mathsf L\mathsf D\mathsf E>0$,
and its constant coefficient is positive. Thus the coefficient signs of
$p(-x)$, in descending powers, are $(+,+,\text{arbitrary},-,+)$.
Descartes' rule of signs gives exactly two sign changes, regardless of
the quadratic coefficient. The two negative roots are therefore real,
distinct, and simple. Write them as $\ell_1<\mu<\ell_2<0$.
The remaining two roots have positive product and positive sum, because
$p(0)>0$ and the sum of all four roots is $\mathsf A+d>0$.
If real, both are positive; if complex, their real parts are positive.
This proves \eqref{eq:rewrite-uncapped-unit-local-spectrum}.

The eigenvalue count does not by itself ensure that a convergent path can
be selected for every nearby pair of initial stocks. I establish this by
showing that the stock components of the stable eigenvectors are linearly
independent. Write a stable eigenvector in transformed coordinates as
$(x,y,z,u)'$.
The component $y$ cannot be zero: the second, fourth, and first rows
of the eigenvalue equation for $\widehat J$ would successively imply $u=x=z=0$.
Normalize $y=1$. The second and fourth rows then give
\[
 x(\ell)=\frac{\mathsf D\mathsf F+d\ell-\ell^2}{\mathsf D\mathsf E}.
\]
The ordering of the roots implies $x(\ell_1)<0<x(\ell_2)$.
Hence the two projected stable eigenvectors $(x(\ell_1),1)'$ and
$(x(\ell_2),1)'$ are linearly independent. Returning to $(\xi_K,\xi_B)$
only rescales their second component by $1/\nu>0$, so
\eqref{eq:rewrite-uncapped-unit-local-projection} follows as well.

These matrix properties now provide a selection rule for the nonlinear
system. The exact normalized vector field is continuously differentiable
near $\xi=0$. The spectral property just proved and the stable-manifold
theorem supply a two-dimensional local manifold
whose trajectories remain near zero and converge exponentially.
The derivative of its projection onto $(\xi_K,\xi_B)$ is $V_{KB}$.
The nonsingularity just proved and the
inverse-function theorem therefore make this manifold locally the
graph of a continuously differentiable function selecting $(\xi_C,\xi_q)$
from the two initial stock deviations. Returning from log deviations to
levels is a smooth invertible change of variables, so the same holds for
$(C_0,q_0)$ as functions of $(K_0,B_0)$.

\par\medskip\noindent\textit{Step 4. Reconstruction of nearby equilibria.}

I first choose a neighborhood of initial states so that the selected
trajectories stay within the region where the local construction is valid.
For some constants $a,H>0$, trajectories on a sufficiently small stable
manifold satisfy
\[
 \|\xi(t)\|\leq H e^{-at}\|\xi(0)\|,\qquad t\geq0.
\]
For any sufficiently small open neighborhood $\mathcal V$ of zero,
choose a ball contained in $\mathcal V$. Continuity of the graph selecting
the two jump variables permits an open neighborhood $\mathcal U$ of the
initial stock deviations for which $H\|\xi(0)\|$ is smaller than that
ball's radius. Every initial stock pair in $\mathcal U$ therefore has a
trajectory staying in $\mathcal V$ for all $t\geq0$ and converging to zero.
Exponentiation maps $\mathcal U$ to an open neighborhood of
$(K_0^*,B_0^*)$. Shrinking these neighborhoods preserves positivity,
feasibility, and existence for every $t\geq0$. Local uniqueness follows
because every trajectory remaining in this local region and converging
to zero belongs to the same stable manifold, which is a graph over stocks.
This does not assert convergence from arbitrary initial $C_0,q_0$ or
exclude equilibrium trajectories that leave the region.

The statement also permits nearby $A_0,N_0$. Equations
\eqref{eq:rewrite-uncapped-unit-initial-b}
and~\eqref{eq:rewrite-uncapped-unit-initial-y-k} make the reference
$B_0^*,K_0^*$ smooth functions of these two positive initial values.
After normalization by the corresponding BGP, the exact vector field
depends only on the parameters and BGP ratios, not on $A_0,N_0$.
The set of $(A_0,N_0,K_0,B_0)$ whose stock deviations belong to
$\mathcal U$ is therefore open and contains the reference initial condition.
From each such initial condition, the construction converges to the BGP
associated with its own $A_0,N_0$, not necessarily to the same levels as
the original reference path.

Since $f(\xi(t))\to0$, the growth rates of $K,B,C,q$ converge to their
BGP values, not only normalized levels. The static equations give the same convergence for
the interest rate, wage growth, and allocation ratios. The unit-elastic
labor income share remains $(1-\alpha)\omega_L$ at every date.

Exponential convergence implies that $r_t-r^*$ is integrable. Thus
$D_t/e^{-r^*t}$ tends to a finite positive constant, and both TVC expressions
equal their BGP counterparts in
\eqref{eq:rewrite-uncapped-unit-household-tvc}
and~\eqref{eq:rewrite-uncapped-unit-developer-tvc} times factors tending
to finite positive limits. Both TVCs hold. Output and all expenditure
flows are bounded multiples of $e^{g_Y^*t}$, while log consumption per
person is a linear function of time plus a bounded deviation. Both
objective integrals therefore remain finite. The global concavity
verification established above
applies at these new price and allocation paths, proving agent optimality.
Static optimality and market clearing then establish equilibrium, not
merely a solution of the necessary conditions.
\end{proof}

% Additional BGP comparisons retained in the source, outside the proof.
%The BGP allocations also give the income shares
%\begin{equation}
% \frac{wL}{Y}=(1-\alpha)\omega_L>0,\qquad
% \frac{(r+\delta)K}{Y}=\alpha,\qquad
% \frac{\Pi}{Y}=(1-\alpha)\omega_X-u^*-m^*.
% \label{eq:rewrite-uncapped-unit-distribution}
%\end{equation}
%Within the restrictions of Proposition~\ref{prop:rewrite-uncapped-unit-bgp},
%increasing the AI weight raises the common growth rate and the interest rate:
%\begin{equation}
% \frac{\partial g_{Y/N}^*}{\partial\omega_X}
% =\frac{\partial g_w^*}{\partial\omega_X}
% =\frac{\partial r^*}{\partial\omega_X}
% =\frac{\eta(1-\eta)(n+\gamma)}{(1-\eta-\omega_X)^2}>0.
% \label{eq:rewrite-uncapped-unit-comparative-statics}
%\end{equation}
%Labor's income share falls across these parameterizations, but remains
%constant over time on each BGP. As $\omega_X\to0^+$, the common per-worker
%growth rate converges to $\gamma$ and $r^*\to\rho+\gamma$.
%These comparisons concern BGPs, not a transition from a no-AI economy.


\Needspace{6\baselineskip}
\begin{proof}[Proof of Proposition~\ref{prop:rewrite-research-scale}]
\label{proof:rewrite-research-scale}
I construct a sequence of research trajectories whose discounted net profit
grows without bound, although each trajectory has finite expenditure and
AI efficiency. Write $\theta=(1-\alpha)/\alpha$. For a research trajectory
$\{M_t\}_{t\in[0,T]}$, discounted net profit on the fixed horizon is
\[
 J_T(M)=\int_0^T
 \exp\left\{-\int_0^t r_s\dd s\right\}
 [\pi(B_t)-M_t]\dd t.
\]
Integrability of $r$ bounds the discount factor above and away from zero
on $[0,T]$.
At the optimal service choice, equations~\eqref{eq:rewrite-ai-price}
and~\eqref{eq:rewrite-monopoly-foc} imply
\[
 \pi(B)=p_XX-\frac XB
 =e_Xp_XX=(1-\alpha)e_Xs_XY.
\]
For $\sigma>1$, marginal revenue is positive at every finite $X$
and approaches zero as $X\to\infty$, as established in
Lemma~\ref{lem:rewrite-monopoly-choice}. The monopoly condition
therefore implies $X\to\infty$ as $B\to\infty$, holding $K,A,L$
fixed. Consequently, $s_X\to1$ and $e_X\to\alpha$.
Substituting these limits into
equation~\eqref{eq:rewrite-substitutes-dated-output-capital} gives
\[
 \pi(B)\sim\kappa_XKB^\theta,\qquad
 \kappa_X=\alpha(1-\alpha)
 \left[(1-\alpha)^2\omega_X^{\sigma/(\sigma-1)}\right]^\theta>0.
\]
The operating-profit asymptotic is uniform over dates in $[0,T]$ because
$K,A,L$ are continuous, bounded, and bounded away from zero
on $[0,T]$. Thus a common lower bound proportional to $B^\theta$ applies
at every date once $B$ is sufficiently large.

The research trajectories constructed below are deviations in the
developer's optimization problem, not alternative equilibrium allocations.
Fix $\Delta\in(0,T)$, choose $M=\mathcal M$ on $[0,\Delta]$, and choose
$M=0$ afterward. Integrating the uncapped RSI technology
\eqref{eq:rewrite-uncapped-unit-law} over the research burst gives
\[
 B_{\mathcal M}(t)^{1-\eta}=B_0^{1-\eta}
 +(1-\eta)\chi\mathcal M^\eta\min\{t,\Delta\}.
\]
As $\mathcal M\to\infty$, efficiency after the research burst is
asymptotically proportional to
$\mathcal M^{\eta/(1-\eta)}$. Discounted operating profit over the remaining
interval therefore grows in proportion to
$\mathcal M^{\eta\theta/(1-\eta)}$, while research expenditure grows only
linearly in $\mathcal M$.
Since optimized operating profit is nonnegative, a lower bound on discounted
net profit counts only operating profits earned after research stops, while
subtracting all research expenditure. For all sufficiently large $\mathcal M$,
\begin{equation}
 J_T(\mathcal M)\geq c_\pi\mathcal M^{
 \eta(1-\alpha)/[\alpha(1-\eta)]}
 -c_M\mathcal M,
 \label{eq:rewrite-research-burst-payoff}
\end{equation}
where $c_\pi,c_M>0$ are independent of $\mathcal M$.
The exponent exceeds one because $\eta>\alpha$,
so $J_T(\mathcal M)\to+\infty$ as $\mathcal M\to\infty$.
Each finite $\mathcal M$ keeps AI efficiency and discounted net profit
finite on $[0,T]$; it is the sequence of payoffs that is unbounded.

Now consider an infinite-horizon equilibrium candidate
$\{\widehat B_t,\widehat M_t\}$ with finite discounted net profit
$\widehat V_0$, and let $\widehat J_T$ denote its payoff through $T$.
After $T$, continue the deviation using $\widehat M_t$.
For sufficiently large $\mathcal M$, the deviation has
$B_{\mathcal M}(T)>\widehat B_T$. Because
$\mathrm d(B^{1-\eta})/\mathrm dt=(1-\eta)\chi M^\eta$,
this efficiency advantage persists under the common subsequent expenditure.
Operating profit is increasing in efficiency by
equation~\eqref{eq:rewrite-profit-envelope}, so the continuation payoff is
at least $\widehat V_0-\widehat J_T$. The full deviation therefore yields
\[
 J_\infty(\mathcal M)\geq J_T(\mathcal M)
 +\widehat V_0-\widehat J_T\longrightarrow+\infty.
\]
No candidate with finite discounted net profit can be optimal.
Moreover, every equilibrium must have finite discounted developer profit.
Indeed, multiplying the household budget~\eqref{eq:rewrite-household-budget}
by the discount factor and integrating, using the Euler equation~\eqref{eq:rewrite-euler}
and transversality condition~\eqref{eq:rewrite-household-tvc} gives
\[
 \int_0^\infty
 e^{-\int_0^t r_s\dd s}(w_tN_t+\Pi_t)\dd t
 =\frac{C_0}{\rho-n}-K_0.
\]
Since labor income is nonnegative, discounted developer profit is
either finite or minus infinity. A payoff of minus infinity cannot be
optimal because
stopping research yields nonnegative operating profit.
Hence no equilibrium exists under the stated restrictions.
This argument does not establish a finite-time singularity.
\end{proof}
